\section{Polar curve splits into conic and straight line}\label{split}

First we describe the types then we prove that the list is complete.
Some types are one-parametric families and we denote the parameter value by $c$. The polar conic will be given either by an explicit parametrization of the circle equations and we reserve $u$ for the parameter, or by indicating the conic equations. The former representation gives also a parametrization of the polar conic: to a circle
$
\epsilon \bigl(x^2+y^2\bigr)+\alpha x+\beta y+\gamma=0$
(where $\epsilon=0$ or $\epsilon=1$, the case $\epsilon=0$ giving a line)
corresponds the polar point with the tetracyclic coordinates
$
[\alpha:\beta:\gamma-\epsilon:-\gamma-\epsilon]$.
The parameter for circles in the pencil will be denoted by~$v$.

To check hexagonality, one computes the Blaschke curvature using the formula in Appendix~\ref{AppendixA} as follows. The polar conic gives a one-parameter family of circles on the unit sphere.
The stereographic projection from the ``south'' pole
\[
X=\frac{2x}{1+x^2+y^2}, \qquad Y=\frac{2y}{1+x^2+y^2}, \qquad Z=\frac{1-x^2-y^2}{1+x^2+y^2},
\]
 transforms this family into a family of circles in the plane parametrized by the points of the conic. The ODE for circles, defined by the conic,
\[
P^2+A(x,y)P+B(x,y)=0, \qquad P=\frac{{\rm d}y}{{\rm d}x},
\]
 is obtained by differentiating and excluding the coordinates of the conic points.
 The slope $R$ comes from the pencil of circles.

For some types, we add an additional geometric detail in the ``title'' to separate the types.

 {\bf 1. Polar conic plane does not cut Darboux quadric, hyperbolic pencil.}
 The pencil with limit circles at the origin and in the infinite point gives circles $x^2+y^2=v$, the polar conic is the circle
$
 X_0^2 + Y_0^2+4cX_0Z_0 =0$, $ U_0=0$, $ c> 0$,
 defining the family
 \[
 x^2+ y^2 -\frac{4c}{c^2u^2 + 1}x +\frac{4c^2u}{c^2u^2 + 1}y=1,
 \]
 the circles of the family enveloping the cyclic
 \[ \bigl(x^2 + y^2\bigr)^2 - \bigl(x^2 + y^2\bigr)(4cx + 2) -4c^2y^2 + 4cx + 1=0
 \]
 as shown on the left in Figure~\ref{123}.


\begin{figure}[t]\centering
\includegraphics[width=0.30\textwidth, trim={6.1cm 2cm 0.2cm 2cm}, clip]{Figure-05-1}
\includegraphics[width=0.33\textwidth, trim={0.3cm 0.3cm 0.3cm 0.3cm}, clip]{Figure-05-2}
\includegraphics[width=0.30\textwidth, trim={0.3cm 3.5cm 8cm 3.5cm}, clip]{Figure-05-3}
 \caption{Types 1, 2, 3.}\label{123}
\end{figure}




 {\bf 2. Polar conic plane does not cut Darboux quadric, hyperbolic pencil, webs symmetric by rotations.}
 The pencil with limit circles at the origin and at the infinite point gives circles $x^2+y^2=v$, the polar conic is the circle
\smash{$
 X_0^2 + Y_0^2 = \frac{4Z_0^2}{c^2}$}, $ U_0=0$,
 defining the family
 \[
 x^2+ y^2 + \frac{2\cos(u)}{c}x+\frac{2\sin(u)}{c} y=1,
 \]
 the circles of the family enveloping the cyclic{\samepage \[
 c^2\bigl(x^2 + y^2\bigr)^2 - \bigl(2c^2 + 4\bigr)\bigl(x^2 + y^2\bigr) + c^2=0,
 \]
 which splits into two concentric circles as shown in the center of Figure~\ref{123}.}

{\bf 3. Polar conic plane cuts Darboux quadric, hyperbolic pencil.}
 The pencil with limit circles at the origin and at the infinite point gives circles
$
x^2+y^2=v$,
the polar conic is the circle
$
x_0^2 + y_0^2=4cx_0$, $ z_0=0$, $ c>0$,
 defining the family
 \[
 x^2+ y^2 +\frac{4c}{c^2u^2 + 1} x - \frac{4c^2u}{c^2u^2 + 1}y=-1,
 \]
 the circles of the family enveloping the cyclic \[
 \bigl(x^2 + y^2\bigr)^2 + \bigl(x^2 + y^2\bigr)(4cx + 2) - 4c^2y^2 + 4cx + 1=0,
 \]
 as shown on the right of Figure~\ref{123}.

\begin{figure}[t]\centering
\includegraphics[width=0.3\textwidth, trim={0.4cm 0.3cm 0.3cm 0.3cm}, clip]{Figure-06-1}
\includegraphics[angle=0, width=0.33\textwidth, trim={0.5cm 3.0cm 0.0cm 0.5cm}, clip]{Figure-06-2}
\includegraphics[angle=0, width=0.3\textwidth, trim={0.4cm 0.0cm 0.0cm 0.0cm}, clip]{Figure-06-3}
 \caption{Types 4, 5, 6.}\label{456}
\end{figure}


 {\bf 4. Polar conic plane cuts Darboux quadric, hyperbolic pencil, webs symmetric by rotations.}
 The pencil with limit circles at the origin and at the infinite point gives circles~${
x^2+y^2=v}$,
the polar conic is the circle
$
x_0^2 + y_0^2=\frac{4}{c^2}$, $ z_0=0$,
 defining the family
 \[
 x^2+ y^2 -\frac{2\cos(u)}{c}x-\frac{2\sin(u)}{c} y=-1,
 \]
 the circles of the family enveloping the cyclic \[
 c^2\bigl(x^2 + y^2\bigr)^2 + \bigl(2c^2 - 4\bigr)\bigl(x^2 + y^2\bigr) + c^2=0,\]
 which splits into two concentric circles as shown on the left in Figure~\ref{456}.



 {\bf 5. Polar conic plane cuts Darboux quadric, elliptic pencil, webs symmetric by homotheties.}
 The pencil with vertexes at the origin and at the infinite point gives lines
$
y=vx$,
the polar conic is
$
y_0^2+ cz_0^2= c$, $ x_0=0$, $ c\notin [0,1]$,
defining the family of circles
$
x^2 + (y + u)^2 =(1-1/c)u^2$,
the circles enveloping the lines $x^2=(c-1)y^2$,
real for $c>1$, as shown in the center in Figure~\ref{456}. The webs of the family are symmetric by homotheties with the center in the origin.


 {\bf 6. Polar conic plane cuts Darboux quadric, elliptic pencil, polar conic and dual to polar line intersect in 2 points on the Darboux quadric.}
 The pencil with vertexes at the origin and at the infinite point gives lines $y=vx$, the polar conic is
$
2cy_0=z_0^2-1$, $ x_0=0$, $ c>0$,
 defining the family of circles
\[
x^2 + y^2 +\frac{(u - 1)}{cu} y + \frac{u - 1}{u + 1}=0,
\]
the circles enveloping the cyclic
\[
c^2\bigl(x^2 + y^2\bigr)^2 + \bigl( 4cy-2c^2 \bigr)\bigl(x^2 + y^2\bigr) + (2y+c)^2=0\]
 as shown on the right in Figure~\ref{456}.

 \begin{figure}[t]\centering
\includegraphics[width=1\textwidth]{Figure-07}
 \caption{Types 7, 8, 9.}\label{789}
\end{figure}


 {\bf 7. Polar conic plane cuts Darboux quadric, elliptic pencil.}
 The pencil with vertexes at the origin and at the infinite point gives lines
$y=vx$,
 the polar conic is
$
 2cy_0z_0+1=z_0^2$, $ x_0=0$, $ c>0$,
 defining the family of circles
\[
x^2 + y^2 - \frac{c+u}{cu}y + \frac{c + u}{c - u}=0,
\]
 as shown on the left in Figure~\ref{789}.


 {\bf 8. Polar conic plane cuts Darboux quadric, parabolic pencil.}
 The pencil with the vertex at the infinite point gives lines $y=v$, the polar conic
$
x_0^2 - z_0 =1$, $ y_0=0$,
 gives the family of circles
\smash{$
\big(x - \frac{1}{u}\big)^2 + y^2 = \frac{u^2 - 1}{u^2}$}.
The circles envelopes the ellipse $x^2 + 2y^2 = 2$ as shown in the center in Figure~\ref{789}.


 {\bf 9. Polar conic plane cuts Darboux quadric, parabolic pencil, webs symmetric by translations.}
 The pencil with the vertex at the infinite point gives lines $y=v$, the polar conic~is
$
x_0^2 +\bigl(1-\sqrt{2}\bigr)z_0^2 - 2\sqrt{2}z_0 = 1+\sqrt{2}$, $ y_0=0$,
defining the family of circles
$
(x +u)^2 + y^2 =1$,
the circles touching the lines $y=\pm 1$ as shown on the right in Figure~\ref{789}.




 {\bf 10. Polar conic plane tangent to Darboux quadric, hyperbolic pencil.}
 The pencil with limit circles at $(1,0)$ and $(-1,0)$ gives circles
$
x^2 + y^2+ 1=vx$,
 the polar conic is
$
x_0^2+ {(1-c)y_0^2} =c$, $ z_0=-1$, $ c>0$, $ c\ne 1$,
 defining the family of lines $x_0x + y_0y = 1$, the lines of the family enveloping the conic
$
 cx^2 + \frac{c}{c - 1}y^2=1
$
 with foci $(1,0)$ and $(-1,0)$ as shown on the left in Figure~\ref{101112}.


\begin{figure}[t]\centering
\includegraphics[width=0.30\textwidth, trim={9.2cm 6.35cm 2.8cm 6.4cm}, clip]{Figure-08-1}
\includegraphics[width=0.30\textwidth, trim={3.1cm 3.5cm 3.5cm 3.5cm}, clip]{Figure-08-2}
\includegraphics[width=0.36\textwidth, trim={3.3cm 6cm 3.5cm 4cm}, clip]{Figure-08-3}
 \caption{Type 10, 11, 12.}\label{101112}
\end{figure}


 {\bf 11. Polar conic plane tangent to Darboux quadric, hyperbolic pencil, polar line meets polar conic.}
 The pencil with limit circles at $(1,0)$ and $(-1,0)$ gives circles
$
x^2 + y^2+ 1=vx$,
the polar conic is
$
y_0^2 -cx_0y_0 + cy_0 + x_0=0$, $ z_0=-1$, $ c\ge 0$,
 defining the family of lines~${
x_0x + y_0y = 1}$,
the lines of the family enveloping the parabola{\samepage
\[
(cx - y)^2 - (2c^2 + 4)x - 2cy + c^2=0
\] with focus at $(1,0)$ as shown in the center in Figure~\ref{101112}.}


 {\bf 12. Polar conic plane tangent to Darboux quadric, hyperbolic pencil, polar line contains the point where polar conic plane touches the Darboux quadric.}
 The pencil with limit circles at the origin and infinity gives circles
$
x^2 + y^2=v$,
 the polar conic is
$
x_0^2 + y_0^2 - 2y_0=0$, $ z_0=-1$
defining the family of lines
$
x_0x + y_0y = 1$,
the lines of the family enveloping the parabola $ x^2 + 2y = 1$ with focus at $(0,0)$ as shown on the right in Figure~\ref{101112}.


 {\bf 13. Polar conic plane tangent to Darboux quadric, hyperbolic pencil, web symmetric by rotations.}
 The pencil with limit circles at the origin and infinity gives circles $x^2 + y^2=v$, the polar conic is $x_0^2 + y_0^2 =1$, $ z_0=-1$ defining the family of lines $x_0x + y_0y = 1$, the lines of the family enveloping the circle $ x^2 + y^2 = 1$ as shown on the left in Figure~\ref{131415}.


\begin{figure}[t]\centering
\includegraphics[width=0.315\textwidth, trim={0.5cm 0.5cm 0.5cm 0.5cm}, clip]{Figure-09-1}
\includegraphics[width=0.335\textwidth, trim={5cm 3.2cm 5cm 7.8cm}, clip]{Figure-09-2}
\includegraphics[width=0.325\textwidth, trim={5.5cm 5cm 5.5cm 6.5cm}, clip]{Figure-09-3}
 \caption{Type 13, 14, 15.} \label{131415}
\end{figure}



 {\bf 14. Polar conic plane tangent to Darboux quadric, elliptic pencil.}
 The pencil with vertexes at $(1,0)$ and $(-1,0)$ gives circles
\[
\frac{\bigl(x^2 + y^2-1\bigr)^2}{\bigl(x^2 + y^2-2x+1\bigr)\bigl(x^2 + y^2+2x+1\bigr)}=v,
\]
the polar conic is
$
cx_0^2+ (c + 1)y_0^2 + 1=0$, $ z_0=-1$,
 defining the family of lines
$
x_0x + y_0y = 1$,
 the lines of the family enveloping the conic
\smash{$
 \frac{x^2}{c} + \frac{y^2}{c + 1}=-1
 $} with foci $(1,0)$ and $(-1,0)$ as shown in the center in Figure~\ref{131415}.

 {\bf 15. Polar conic plane tangent to Darboux quadric, parabolic pencil.}
 The pencil with vertex at the origin gives circles
$
x^2 + y^2=2vy$,
the polar conic is
$
x_0^2 + y_0^2 =1$, $ z_0=-1
$ defining the family of lines $x_0x + y_0y = 1$, the lines of the family enveloping the circle $ x^2 + y^2 = 1$ as shown on the right in Figure~\ref{131415}.

\begin{Theorem}
The webs of different types in the above classification list are not M\"obius equivalent. The webs of the same type with different normal forms are not M\"obius equivalent.
\end{Theorem}
\begin{proof} The webs from different types are not M\"obius equivalent: discrete geometric invariants indicated in the descriptions, such as 1) presence of infinitesimal symmetry and 2) the mutual position of polar line, polar conic and Darboux quadric, effectively separate the types.

To see that the different normal forms within a family are not M\"obius equivalent, one computes the subgroup $G_p$ of ${\rm PSO}(3,1)$ respecting the positions of the polar conic plane and the polar line in the chosen normalization and checks that the $G_p$-orbits of the canonical forms are different.

For the first 4 types, $G_p$ is generated by $R_z$ and by reflections in the coordinate planes.

For the 5th, 6th, 7th type, $G_p$ is generated by $B_z$ and by reflections in the coordinate planes.

For the 10th, 11th and 14th types, $G_p$ is discrete and generated by reflections in the planes $x=0$ and $y=0$.
\end{proof}


As in the case of 3 pencils, Lemma \ref{sing} effectively selects candidates among 3-webs that can be hexagonal. The singularities of the type described by the Lemma arise when either 1) a line, joining 2 different points on the polar conic, touches the Darboux quadric at a point $p_t$, while the polar conic plane is not tangent the Darboux quadric or 2) a line, joining a point on the polar line with a point on the polar conic, touches the Darboux quadric at a point $p_t$ while the tangent plane to the Darboux quadric at $p_t$ is not tangent to the polar conic.

In the former case, the points $p_t$ trace a circle, which is the intersection of the polar conic plane with the Darboux quadric. Then, by Lemma \ref{sing}, the polar of the polar conic plane lies on the pencil polar line.

In the latter case, consider a point $p_r$ running over the polar line. For non-planar polar set, $p_r$ meets the polar conic plane $\pi_c$ only at one point. Therefore, the one-parameter family of cones tangent to the Darboux quadric and having their vertexes at $p_r$ cuts the plane $\pi_c$ in a one-parameter family of conics $c_r$. Each conic $c_r$ intersect the polar conic $c_p$ at 4 (possibly complex or multiple) points. If the polar conic $c_p$ is not a member of the family $\{c_r\}$, at least one of these 4 intersection points is moving along $c_p$ as $p_r$ runs over the polar line. In fact, if intersection points are stable then $\{c_r\}$ is a pencil of conics containing $c_p$.

Choose one such moving intersection point $p_i$. Lines $l_r$ tangent to $c_r$ at $p_i$ form a one-parameter family, or congruence of lines. All the objects in this construction are considered as complex but the polar conic and the polar line must have equations with real coefficients and the real part of the polar conic cannot lie completely inside the Darboux quadric.

\begin{Proposition}\label{geometryconic}
If the polar curve of a hexagonal circular $3$-web is non-planar and splits into a line and a smooth conic, then for the web complexification hold true
\begin{itemize}\itemsep=0pt
\item[$(1)$] the polar of the polar conic plane lies on the polar line, if this plane is not tangent to the Darboux quadric and
\item[$(2)$] the congruence of lines $l_r$ is a pencil with the vertex on the polar conic, if the polar conic is not a member of the family $\{c_r\}$.
\end{itemize}
\end{Proposition}
\begin{proof}
 The first claim follows directly from Lemma \ref{sing}: the points, where bisecant lines touch the Darboux quadric, trace a circle whose polar point must lie on the polar line. For conic planes missing the Darboux quadric, the complex version works.

To derive the second claim, observe that the line $p_rp_i$ touches the Darboux quadric at a~singular point treated by Lemma \ref{sing}. Thus the tangency point must trace a circle on the Darboux quadric and the polar point $p_0$ of the circle must lie on the polar conic. Then the plane, tangent to the Darboux quadric and passing through $p_rp_i$, contains $p_0$. This plane cuts the plane $\pi_c$ of the polar conic along a line $l_r$, passing through the $p_i$ and tangent to the corresponding conic~$c_r$. Hence $\{l_r\}$ is the pencil with vertex at $p_0$.
\end{proof}

\begin{Theorem}\label{conicPencil}
If the polar curve of a hexagonal circular $3$-web splits into a smooth conic and a straight line, not lying in a plane of the conic, then the web is M\"obius equivalent to one from the above presented list of types $1$--$15$.
\end{Theorem}
\begin{proof} The polar conic plane either completely misses the Darboux quadric, or cuts it in a~circle, or is tangent to it. The polar line is either hyperbolic, or elliptic, or parabolic. Thus we have 9 cases to consider, each case defining a set of webs (possibly empty).

$\bullet$ {\it Polar conic plane misses the Darboux quadric}.
Applying a suitable M\" obius transformation, we can send the polar conic plane to infinity. Then by Proposition \ref{geometryconic} the polar line contains the origin of the affine chart and therefore meets the Darboux quadric at 2 points. Thus the polar line is hyperbolic. Applying a rotation around the origin, we map these two points to~${(0,0,\pm 1)}$. In the affine coordinates $x=\frac{X}{Z}$, $ y=\frac{Y}{Z}$ on the polar conic plane $U=0$, the conics~$c_r$ are~${
x^2+y^2=r}$,
where $r$ is considered as a complex parameter. These conics are real only for real non-negative $r$.
If the polar conic coincides with one of the conics~$c_r$, then we get Type~2, symmetric by~$R_z$.

If the polar conic is not one of $c_r$, then, by Proposition \ref{geometryconic}, the points, where lines of pencil with the vertex at some point $p_0=(x_0,y_0)\in \mathbb C^2$ touch the circles $c_r$, run over the polar conic. The line of the pencil $y=y_0+k(x-x_0)$ corresponding to the parameter $k\in \mathbb C$, is tangent to the conic $c_r$ for
\[
r=\frac{x_0^2k^2 - 2x_0y_0k + y_0^2}{k^2 + 1}
\]
at the point
\[
p(k)=(x(k),y(k))=\left(\frac{k(x_0k - y_0)}{k^2 + 1},\frac{y_0-x_0k}{k^2 + 1} \right).
\]
The points $p(k)$ run over the complex conic
$
x^2+y^2 =x_0x + y_0x$.
 This conic is real if and only if~$x_0$ and $y_0$ are real. It is smooth if and only if $p_0\ne (0,0)$. We got the first web of the list, Type~1. This conic is the circle, passing through the origin $O=(0,0)$ and $p_0$ and having its center at the midpoint of the segment $Op_0$.
(We obtained a theorem of scholar geometry.) Using~$R_z$, we normalize $p_0$ to $y_0=0$, $x_0>0$.

 For infinite $p_0$, different from the cyclic points, the line $l(k)$ of the pencil $y=ax+k$, $a\in \mathbb C$ touches a unique conic of the family $\{c_r\}$ at the point
$
p(k)=\bigl(-\frac{ak}{a^2 + 1}, \frac{k}{a^2 + 1}\bigr)$.
The points $p(k)$ are collinear: $x(k)+ay(k)=0$ and we cannot obtain a smooth polar conic in this way. Finally, if $p_0$ is cyclic, the lines from the pencil touch the conics $c_r$ at the very point~$p_0$ and we get no conic at all.


$\bullet$ {\it Polar conic plane cuts the Darboux quadric, hyperbolic pencil}.
We send the conic plane to~${Z=0}$. Now the polar line contains the infinite point $[0:0:1:0]$ by Proposition \ref{geometryconic}. The subgroup of the M\" obius group, preserving the plane $Z=0$, is generated by rotations around $Z$-axis and the boosts $B_x$ and $B_y$. Using this subgroup, we normalize the polar line to~${x=y=0}$. The conics $c_r$ are circles in the plane $z=0$ with the center at the origin. Repeating the arguments that we used above for conic planes missing the Darboux quadric, we get Types 3 and 4.

$\bullet$ {\it Polar conic plane cuts the Darboux quadric, elliptic pencil}.
We normalize the conic plane to $X=0$. Then by Proposition \ref{geometryconic} the point $[1:0:0:0]$ lies on the polar line. The polar line, being elliptic, meets the plane $X=0$ at a point $p$ outside the unit circle. M\" obius transformations, preserving the plane $X=0$, are generated by rotations around $X$-axis and the boosts $B_y$ and $B_z$. Using these transformations, we send the point $p$ to $[0:1:0:0:0]$. Now the polar line is $U=Z=0$ and the conics $c_r$ have equations
\smash{$
z^2+\frac{ y^2}{r}=1
$} in the affine coordinates.
If the polar conic coincides with one of these conics, then $r$ is real and we get Type 5.


 Otherwise, by the second claim of Proposition \ref{geometryconic}, the lines $l_r$ meet at some point $p_0\in c_p$. If $p_0=(0,y_0,z_0)$ is finite, a line of the pencil of lines $z=z_0+k(y - y_0)$ with vertex at $p_0$ is tangent to the conic $c_r$ with
\[
r=\frac{y_0^2k^2 - 2y_0z_0k + z_0^2 - 1}{k^2}.
\]
 Thus the points, where the lines $l_r$ touch $c_r$, are parametrized by $k$ via
\[
y=\frac{y_0^2k^2 - 2y_0z_0k + z_0^2 - 1}{k(y_0k - z_0)}, \qquad z=\frac{1}{z_0-y_0k}.
\]
Excluding $k$, we get the conic
$
y_0z^2-z_0yz + y - y_0=0$.
This conic is real only for real $y_0$, $z_0$ and is smooth if and only if $y_0\bigl(z^2_0-1\bigr)\ne 0$.
 If $z^2_0<1$, we normalize to $z_0=0$ applying $B_z$ and get Type~6. If $z^2_0> 1$, we send $p_0$ to infinity applying~$B_z$.

For infinite point $p_0=[0:Y_0:Z_0:0]$, we conclude that $Y_0\ne 0$ and $Z_0\ne 0$, otherwise the tangency points $p_i$ do not trace a conic. Thus we can set $p_0=[0:y_0:1:0]$, where $y_0\ne 0$. In~the affine coordinates $u=\frac{U}{Z}$, $y=\frac{Y}{Z}$ in the plane $X=0$, the conic $c_r$ is \smash{$1+\frac{y^2}{r}=u^2$}. A~line from the pencil $y=ku+y_0$ is tangent to the conic $c_r$ if and only if
$
r=k^2 - y_0^2$.
The points, where the lines $l_r$ touch $c_r$, are parameterized by $k$ via
\smash{$
y=\frac{y_0^2 - k^2}{y_0}$}, \smash{$ u=-\frac{k}{y_0}$}.
This is a parametrization of the polar conic of Type 7
$
y_0u^2 + y = y_0$,
which becomes
$
y_0U^2 - y_0Z^2 + YZ=0
$
in homogeneous coordinates.

$\bullet$ {\it Polar conic plane cuts the Darboux quadric, parabolic pencil}.
We normalize the conic plane to $Y=0$. Then the point $[0:1:0:0]$ is on the polar line.
 The plane $Y=0$ is stable under rotations $R_y$ along the $y$-axis and under boosts $B_x$ and $B_z$.
Rotating, if necessary, around the $y$-axis, we bring the polar line to $z=-1,$ $x=0.$ In the affine coordinates on the plane $Y=0$, the conics $c_r$ have equations
\[
x^2 + \frac{r-1}{r}z^2 - \frac{2}{r}z - \frac{r+1}{r}=0.
\]
If the polar conic coincides with one of $c_r$, then $r$ is real. The chosen position of polar conic plane and polar line is stable by action of the group generated by $B_z$. Applying it, one can normalize to $r=\frac{1}{\sqrt{2}}$ and we get Type 9.

Otherwise, consider first the pencil of lines $z=z_0+k(x - x_0)$ with finite vertex at $p_0=(x_0,0,z_0)$. A line from the pencil is tangent to the conic $c_r$ with
\[
r=\frac{x_0^2k^2 - 2x_0(z_0 + 1)k + (z0 + 1)^2}{\bigl(x_0^2 - 1\bigr)k^2 - 2x_0z_0k + z_0^2 - 1}.
\]
 Thus the points, where the lines $l_r$ touch $c_r$, are parametrized by $k$ via
\[
x=\frac{k(x_0k - z_0 - 1)}{k^2 - x_0k + z_0 + 1}, \qquad z=\frac{\bigl(x_0^2 - 1\bigr)k^2 - x_0(2z_0 + 1)k + z_0(z_0 + 1)}{k^2 - x_0k + z_0 + 1}.
\]
Excluding $k$, we get the conic
\begin{equation}\label{polarcutP}
(z_0 + 1)x^2 - x_0xz +z^2- x_0x+ (1-z_0)z - z_0=0.
\end{equation}

This conic is real only for real $x_0,z_0$ and is smooth if and only if $z_0\ne -1.$
The chosen position of polar conic plane and polar line is stable by action of the group generated by $B_z$ and $B_x-R_y$. Consider the orbits of points in the plane $Y=0.$ The orbit dimension is two for points outside the union of the line $U+Z=0$ and the circle $X^2+Z^2=U^2$, and is one on this union except for their common point.
Thus the finite representatives of the orbits are $[0:0:0:1]$, $[0:0:1:1]$ and $[0:0:-1:1].$ For the first two points, the conics \eqref{polarcutP} lie inside the Darboux quadric and there is no real circles. For the point $[0:0:-1:1]$, the conic \eqref{polarcutP} is not smooth.

For infinite vertexes $p_0$, the pencil $z=k$ gives a non-smooth conic. Thus the pencil can be chosen as $x=az+k$. A line from the pencil is tangent to the conic $c_r$ with
\smash{$
r=\frac{(k-a)^2}{k^2 -a^2 - 1}$}.
 The points, where the lines $l_r$ touch $c_r$, are
\smash{$
x=\frac{k-a}{a^2 - ak + 1}$}, \smash{$ z=\frac{k^2-ak-1}{a^2 - ak + 1}$}.
Excluding $k$, we get the conic~${
x^2-axz-ax-z-1=0}$,
which is real only for real $a$. Taking into account the action of the group generated by $B_z$ and $B_x-R_y$, we set $a=0$ and get Type 8.

$\bullet$ {\it Polar conic plane tangent to Darboux quadric, hyperbolic pencil}.
 If the hyperbolic line does not contain the point where the polar conic plane touches the Darboux quadric, then we sent this point to $(0,0,-1)$ and the polar line to $Y=Z=0$. In the affine coordinates on the plane~${z=-1}$, the conics $c_r$ have equations
$(x - r)^2 + \bigl(1-r^2\bigr)y^2 + 1-r^2=0$.
If the polar conic coincides with one of $c_r$, then the web curvature
\[
K_B=\frac{4\bigl(r^2 - 1\bigr)^2\bigl(x^2 - 1\bigr)\bigl( rx^2 + ry^2 +\bigl(r^2-3\bigr)x +r\bigr)\bigl(x^2 + y^2-2rx + 1\bigr)^4}{x^4y^3\bigl(x^2 +1-2rx \bigr)^6}
\]
vanishes only for $r=\pm 1$, the conic $c_r$ being non-smooth for these values.

A line from the pencil $y=y_0+k(x - x_0)$ with finite vertex at $p_0=(x_0,y_0,-1)$ is tangent to the conic $c_r$ with
\[
r=\frac{\bigl(x_0^2 + 1\bigr)k^2 - 2x_0y_0k + y_0^2 + 1}{2k(x_0k - y_0)}.
\]
 Thus the points, where the lines $l_r$ touch $c_r$, are parametrized by $k$ via
\begin{gather*}
x=\frac{x_0\bigl(x_0^2 - 1\bigr)k^3 - y_0\bigl(3x_0^2 - 1\bigr)k^2 + x_0\bigl(3y_0^2 + 1\bigr)k - y_0\bigl(y_0^2 + 1\bigr)}{k\bigl(\bigl(x_0^2 - 1\bigr)k^2 - 2x_0y_0k + \bigl(y_0^2 - 1\bigr)\bigr)},
\\
y=\frac{2(x_0k - y_0)}{\bigl(x_0^2 - 1\bigr)k^2 - 2x_0y_0k + \bigl(y_0^2 - 1\bigr)}.
\end{gather*}
Excluding $k$, we get the cubic
\begin{equation}\label{polartH}
\bigl(x_0^2 - 1\bigr)y^3 +\bigl(y_0^2 - 1\bigr)x^2y - 2x_0y_0xy^2+ 2y_0x^2+ 2y_0y^2- 2x_0y_0x +\bigl(x_0^2 - y_0^2\bigr)y=0.
\end{equation}
This cubic splits into a smooth real conic and a line in 3 cases:
\begin{itemize}\itemsep=0pt
\item[(1)] for $y_0=0$, \eqref{polartH} factors as $y\bigl(x^2 +\bigl(1-x_0^2\bigr)y^2 - x_0^2\bigr)=0$ and we get Type 10,
\item[(2)] for $x_0= 1$, $y_0\ne 0$, \eqref{polartH} factors as $(x - 1)\bigl(\bigl(y_0^2-1\bigr)xy - 2y_0y^2 +2y_0x+ \bigl(y_0^2-1\bigr)y\bigr)=0$,
\item[(3)] for $x_0= -1$, $y_0\ne 0$, \eqref{polartH} factors as $(x + 1)\bigl(\bigl(y_0^2-1\bigr)xy + 2y_0y^2 +2y_0x- \bigl(y_0^2-1\bigr)y\bigr)=0$.
\end{itemize}
The cases 2) and 3) give Type 11, the substitution $x\to -x$ reducing one to the other.

For infinite vertexes $p_0$, the pencil $x=k$ gives a non-smooth conic. Thus the pencil can be chosen as $y=ax+k$. A line from the pencil is tangent to the conic $c_r$ with
\smash{$
r=\frac{k^2+a^2+1}{2ak}$}.
 The points, where the lines $l_r$ touch $c_r$, are
\[
x=\frac{k\bigl(k^2 -a^2+1\bigr)}{a\bigl(k^2-a^2-1\bigr)}, \qquad y=\frac{2k}{a^2 - k^2 + 1}.
\]
Excluding $k$, we get the cubic
$
y^3+a^2x^2y+2axy^2+2ax+\bigl(1-a^2\bigr)y=0$.
The cubic splits into a line and a conic only for $a=0$ or for $a^2 \pm 2{\rm i}a -1=0$. The former case gives non-real and non-smooth conic $y^2+1=0$, the latter -- the non-real conic $xy \pm {\rm i} \bigl(y^2+2\bigr)=0$.

 If the hyperbolic line contains the point where the polar conic plane touches the Darboux quadric, then we sent this point to $(0,0,-1)$ and the polar line to $X=Y=0$. In the affine coordinates on the plane $z=-1$, the conics $c_r$ are concentric circles. The case of concentric circles was considered above. We get Types 12 and~13.


$\bullet$ {\it Polar conic plane tangent to Darboux quadric, elliptic pencil}.
If one vertex of the elliptic pencil coincides with the point where the polar conic plane touches the Darboux quadric, then the polar curve is planar.

Thus we can sent the tangent point to $(0,0,-1)$ and the vertexes to $(\pm 1,0,0)$. The conics $c_r$ have equations
$
\bigl(r^2 + 1\bigr)x^2 + (y+r)^2 = r^2+1$.
If the polar conic coincides with one of $c_r$, then the web curvature
\[
K_B=-\frac{64\bigl(r^2 + 1\bigr)^2x\bigl(y^2 + 1\bigr)\bigl( rx^2 + ry^2+ \bigl(3+r^2\bigr)y - r \bigr)\bigl(2ry - x^2 - y^2 + 1\bigr)^4}{\bigl(x^2 - y^2 - 1\bigr)^4\bigl(r^2 - x^2 + 1\bigr)^6}
\]
vanishes only for $r=\pm {\rm i}$, the conic $c_r$ being non-smooth for these values.

A line from the pencil $y=y_0+k(x-x_0)$ is tangent to the conic $c_r$ with
\[
r=\frac{\bigl(x_0^2 - 1\bigr)k^2 - 2x_0y_0k + \bigl(y_0^2 - 1\bigr)}{2(x_0k - y_0)}.
\]
 The points, where the lines $l_r$ touch $c_r$, are
\begin{gather*}
x=\frac{2k(x_0k - y_0)}{\bigl(x_0^2 + 1\bigr)k^2 - 2x_0y_0k + y_0^2 + 1},\\
y=-\frac{x_0\bigl(x_0^2 - 1\bigr)k^3 - y_0\bigl(3x_0^2 - 1\bigr)k^2 + x_0\bigl(3y_0^2 + 1\bigr)k - y_0\bigl(y_0^2 + 1\bigr)}{\bigl(x_0^2 + 1\bigr)k^2 - 2x_0y_0k + y_0^2 + 1}.
\end{gather*}
Excluding $k$, we get the cubic
\[
\bigl(y_0^2 + 1\bigr)x^3 - 2x_0y_0x^2y+ \bigl(x_0^2 + 1\bigr)xy^2 - 2x_0x^2 - 2x_0y^2 + \bigl(x_0^2 - y_0^2\bigr)x + 2x_0y_0y=0.
\]
This cubic splits into a conic and a line in 4
 cases:
 \begin{itemize}
\item[(1)] for $y_0=\pm {\rm i}$, the cubic equation factors as
\[
(\pm i-y)\bigl(\pm 2{\rm i}x_0x^2 -\bigl(1+x_0^2\bigr)xy \mp {\rm i}\bigl(1+x_0^2\bigr)x + 2x_0y\bigr)=0,
\]
\item[(2)] for $y_0=\pm {\rm i}x_0$, the cubic equation factors as
\[
 (x \pm {\rm i} y)\bigl(\bigl(1-x_0^2\bigr)x^2 \mp {\rm i} \bigl(x_0^2+1\bigr)xy - 2x_0x \pm2 {\rm i} x_0y + 2x_0^2\bigr)=0,
\]
\item[(3)] for $x_0=0$, the cubic equation factors as $x\bigl(\bigl(y_0^2 + 1\bigr)x^2 + y^2 - y_0^2\bigr)=0$,
\item[(4)] for $x_0=\pm 1$, the cubic equation factors as $(x\mp 1)\bigl(\bigl(y_0^2 + 1\bigr)x^2+2y^2 \mp 2y_0xy \pm \bigl(y_0^2 - 1\bigr)x - 2y_0y\bigr)=0$.
 \end{itemize}
In the cases 1) and 2) the conic is not real, the case 3) gives Type 14, for the case 4) the web curvature does not vanish.



For infinite vertexes $p_0$, the pencil $x=k$ gives a non-smooth conic. Thus the pencil can be chosen as $y=ax+k$. A line from the pencil is tangent to the conic $c_r$ with
\smash{$
r=\frac{a^2 - k^2 + 1}{2k}$}.
 The points, where the lines $l_r$ touch $c_r$, are
\[
x=-\frac{2ak}{k^2+a^2+1}, \qquad y=\frac{k\bigl( k^2-a^2 + 1\bigr)}{k^2+a^2+1}.
\]
Excluding $k$, we get the cubic
\[
a^2x^3 - 2ax^2y + xy^2 +\bigl(1-a^2\bigr)x + 2ay=0.
\]
The cubic splits into a line and a conic in 2 cases:
\begin{itemize}
\item[(1)] for $a=0$ the cubic equation factors as $x\bigl(y^2+1\bigr)=0$ and the conic is non-smooth
\item[(2)] for $a=\pm {\rm i}$ the cubic equation factors as $(x \pm {\rm i} y)\bigl(x^2 \pm {\rm i} xy-2\bigr)=0$, and the conic is not real.
\end{itemize}


$\bullet$ {\it Polar conic plane tangent to Darboux quadric, parabolic pencil}.
For non-planar polar curves, the points, where the polar line and polar conic plane touch the Darboux quadric, are different. Thus we can normalize the polar conic plane to $z=-1$ and the polar line to $x=0$, $z=1$. This configuration is preserved by $B_z$. The conics $c_r$ have equations
$
\frac{r}{4}x^2 + y =\frac{1}{r}$.
If the polar conic coincides with one of $c_r$, then the web curvature
\[
K_B=-\frac{4096r^5xy^2\bigl(ry + 2x^2 + 2y^2\bigr)\bigl(ry - x^2 - y^2\bigr)^4}{\bigl(x^2 - y^2\bigr)^4\bigl(r^2 - 4x^2\bigr)^6}
\]
vanishes only for $r=0$, the conic $c_r$ being non-smooth for this value.

 A line from the pencil $y=y_0+k(x-x_0)$ is tangent to the conic $c_r$ with
$
r=\frac{k^2+1}{y_0-x_0k}$.
 The points, where the lines $l_r$ touch $c_r$, are
\[
x=\frac{2k(x_0k - y_0)}{k^2 + 1},\qquad
y=\frac{x_0k^3 -y_0 k^2 -x_0 k + y_0}{k^2 + 1}.
\]
Excluding $k$, we get the cubic
\[
x^3+ xy^2 - 2x_0x^2- 2x_0y^2 + \bigl(x_0^2 - y_0^2\bigr)x + 2x_0y_0y=0.
\]
This cubic splits into a conic and a line in 2
 cases:
 \begin{itemize}
\item[(1)] for $y_0=\alpha x_0$, where $\alpha ^2 \pm 2{\rm i} \alpha -1=0$, the cubic equation factors as
\[
(x \pm {\rm i} y)\bigl(\pm {\rm i} xy -x^2 +2x_0x\mp 2{\rm i} x_0y - 2x_0^2\bigr)=0,
\]
\item[(2)] for $x_0=0$, the cubic equation factors as $x\bigl(x^2+y^2-y_0^2\bigr)=0$.
 \end{itemize}
 In the case 1), the conic is not real, the case 2) gives Type 15 after rescaling by $B_z$.

 For infinite vertexes $p_0$, the pencil $y=k$ gives a non-smooth conic. Thus the pencil can be chosen as $x=ay+k$. A line from the pencil is tangent to the conic $c_r$ with
$
r=-\frac{a^2+ 1}{ak}$.
 The points, where the lines $l_r$ touch $c_r$, are
\[
x=\frac{2k}{a^2+1}, \qquad y=\frac{k\bigl(1- a^2 \bigr)}{a\bigl(1+a^2\bigr)}.
\]
Excluding $k$, we get the line
$\bigl(a^2 - 1\bigr)x + 2y=0$.

 Thus we have considered all the cases with non-planar polar curve, the theorem is proved.
\end{proof}
